\documentclass[10pt,a4paper]{amsart}
\usepackage[margin=19mm]{geometry}
\usepackage{amsmath,amssymb,mathtools}
\usepackage[scr=rsfs,cal=euler]{mathalfa}
\usepackage{xfrac}
\usepackage[unicode,hidelinks,bookmarks=false]{hyperref}
\newcommand{\ie}{i.e.\ }
\newcommand{\cf}{cf.\ }
\newcommand{\resp}{respectively\ }
\newcommand{\Subsection}{\subsection}
\providecommand{\nobreakdash}{\nobreak\hbox{-}}
\setlength{\emergencystretch}{2em}
\multlinegap0pt
\usepackage{amssymb}
\usepackage{mathtools}
\usepackage[shortlabels]{enumitem}
\usepackage{xcolor}
\newtheorem{definition}{Definition}
\newtheorem{lemma}{Lemma}
\newtheorem{theorem}{Theorem}
\newcommand{\Z}{\mathbb Z}
\newcommand{\one}{\mathbf 1}
\DeclareMathOperator{\row}{row}
\title{Symplectic and projective small covers over products of polygons}
\author{Suyoung Choi}
\date{Source: arXiv:2605.22554v1 (CC BY 4.0)}
\begin{document}
\maketitle

\noindent\emph{Source and licence.} Symplectic and projective small covers over products of polygons. Version arXiv:2605.22554v1 (CC BY 4.0), distributed by arXiv under the Creative Commons Attribution 4.0 International licence, http://creativecommons.org/licenses/by/4.0/, as recorded in the arXiv licence metadata for this exact version and shown on the abstract page https://arxiv.org/abs/2605.22554v1. Full licence notice: \texttt{licenses/arxiv-2605.22554-CC-BY-4.0.txt}.\par\smallskip

\noindent\textit{Editorial scope.} Counted unit: the article's theorem thm:blockize (source0.tex lines 1080--1094). The article's complete original proof of it is delivered with it (source0.tex lines 1095--1170, one \texttt{proof} environment of 76 lines). No mathematics is written, repaired or re-derived: the statement and the proof are the article's own lines, in the article's own notation, and no retained line is substituted or re-flowed.  Retained uncounted as the source-local context the proof needs: lines 460--474; lines 1034--1046.  Nothing was omitted from any retained span, so the package carries no omission record.  No retained line contains a citation, so the wrapper carries no bibliography and no bibliography entry.  No retained line contains a figure environment, an image-inclusion command or drawing code, so no image is delivered and the package declares no asset dependency.  Editorial formatting only, in addition: an AMS \texttt{amsart} preamble that restates the article's own declarations and macros used by the retained lines, namely the environments \texttt{definition}, \texttt{lemma}, \texttt{theorem} on separate counters, together with the standard packages those lines need.  The retained environments print in this document as Definition 1, Lemma 1, Theorem 1 in order of appearance, whereas the article prints the same items with the numbering of its own full document.  The complete native source of the article as distributed in the arXiv e-print for this exact version is bundled unchanged as \texttt{sources/201144/source0.tex}.
\begin{definition}\label{def:factor-compatible}
Let $M=M(P,\lambda)$ be a small cover over $P=\prod_{i=1}^n P_{m_i}$. 
We call $M$ \emph{factor-compatible} if the following condition holds.

After a permutation of the interval directions in the cube part
$$
    \prod_{m_i=4} P_4 \cong (\Delta^1)^{2r}
$$
and after regrouping them into $r$ square factors, the total facet weight $\chi_i$ of every resulting polygon factor satisfies
$$
    \chi_i\in\row(\lambda).
$$
After this regrouping, we relabel the resulting polygon factors and denote their total facet weights again by $\chi_i$.
Whenever a factor-compatible small cover is considered, such a regrouping is fixed once and for all.
\end{definition}

\begin{lemma} \label{lem:colored triangularization}
    Let $E$ be an $n$-dimensional vector space over $\Z_2$, and let
    $$
        D_1,\ldots,D_n\subset E^\ast\setminus\{0\}
    $$
    be nonempty subsets.
    Suppose that every transversal~$d_1\in D_1,\ldots,d_n\in D_n$ is a basis of $E^\ast$.
    Then, after permuting the subsets~$D_i$, there exists a basis $\varepsilon_1,\ldots,\varepsilon_n$ of~$E^\ast$ such that
    $$
        D_i\subset \varepsilon_i+ \langle \varepsilon_{i+1},\ldots,\varepsilon_n\rangle \qquad \text{for every } i=1,\ldots,n.
    $$
    Here the span is interpreted as $0$ when $i=n$.
\end{lemma}

\begin{theorem}\label{thm:blockize}
    Let $P=\prod_{i=1}^n P_{m_i}$ be a product of polygons, and let $\lambda$ be a factor-compatible characteristic matrix on $P$.
    Then, after permuting the polygon factors, $\lambda$ is Davis--Januszkiewicz equivalent to a block lower triangular characteristic matrix
    $$
        \lambda'=
        \begin{pmatrix}
        \Lambda_{11} & 0            & \cdots & 0\\
        \Lambda_{21} & \Lambda_{22} & \cdots & 0\\
        \vdots       & \vdots       & \ddots & \vdots\\
        \Lambda_{n1} & \Lambda_{n2} & \cdots & \Lambda_{nn}
        \end{pmatrix},
    $$
    where each $\Lambda_{ii}$ is a characteristic matrix on $P_{m_i}$.
    Moreover, the first row of each $\Lambda_{ii}$ may be chosen to be the all-one vector.
\end{theorem}

\begin{proof}
    By factor-compatibility, and by the convention fixed in Definition~\ref{def:factor-compatible}, we have $\chi_i\in\row(\lambda)$ for every factor $i$.
    Set $C=\langle \chi_1,\ldots,\chi_n\rangle\subset \row(\lambda)$.
    The supports of the $\chi_i$'s are disjoint, so $\dim C=n$.
    Since $\lambda$ is a characteristic matrix of a $2n$-dimensional small cover, $\dim \row(\lambda)=2n$.
    Hence $E:=\row(\lambda)/C$ has dimension~$n$.
        
    For each factor $i$ and each adjacent pair $F_{i,k},F_{i,k+1}$, define a linear functional $d_{i,k} \colon E\to \Z_2$ by
    $$
        d_{i,k}([w])=w(F_{i,k})+w(F_{i,k+1}).
    $$
    This is well-defined because adding $\chi_i$ changes both values $w(F_{i,k})$ and $w(F_{i,k+1})$ by $1$, while adding $\chi_j$ for $j\neq i$ does not affect these two coordinates.
    Put
    $$
            D_i=\{d_{i,1},\ldots,d_{i,m_i}\}\subset E^\ast.
    $$
        
    We claim that every transversal $d_{1,k_1}\in D_1,\ldots,d_{n,k_n}\in D_n$ is a basis of $E^\ast$.
    Choose the vertex of $P$ corresponding to the adjacent pairs
    $$
        \{F_{1,k_1},F_{1,k_1+1}\},\ldots,\{F_{n,k_n},F_{n,k_n+1}\}.
    $$
    The characteristic condition says that the restriction map $\row(\lambda)\to \Z_2^{2n}$ to these $2n$ facets is an isomorphism.
    Under this restriction, the subspace~$C$ maps to the subspace which is constant on each of the $n$ chosen adjacent pairs.
    Therefore, after quotienting by $C$, the induced map $E \to \Z_2^n$ is given by $[w]\mapsto \bigl(d_{1,k_1}([w]),\ldots,d_{n,k_n}([w])\bigr)$.
    This map is an isomorphism.
    This proves the claim.
    In particular, no element of any $D_i$ is zero.
        
    Apply Lemma~\ref{lem:colored triangularization} to $D_1,\ldots,D_n$.
    After permuting the polygon factors, there exists a basis $\varepsilon_1,\ldots,\varepsilon_n$ of~$E^\ast$ such that $D_i\subset\varepsilon_i+\langle \varepsilon_{i+1},\ldots,\varepsilon_n\rangle$ for all $i$.
    Let $\tau_1,\ldots,\tau_n$ be the dual basis of $E$, and choose arbitrary lifts $\widetilde\tau_i\in \row(\lambda)$ of $\tau_i$.
        
    Fix $i$.
    If $k>i$, then every element of $D_k$ lies in $\varepsilon_k+\langle\varepsilon_{k+1},\ldots,\varepsilon_n\rangle$.
    Hence every $d\in D_k$ satisfies $d(\tau_i)=0$.
    Thus, for every adjacent pair in the $k$-th polygon factor,
    $$
        \widetilde\tau_i(F_{k,j})+\widetilde\tau_i(F_{k,j+1})=0.
    $$
    Since the adjacency graph of a polygon is connected, the restriction of $\widetilde\tau_i$ to $I_k$ is constant.
    Hence it is either $0$ or $\chi_k|_{I_k}$.
    By adding suitable multiples of $\chi_k$, for~$k>i$, we may replace $\widetilde\tau_i$ by $\eta_i=\widetilde\tau_i+\sum_{k>i} c_{ik}\chi_k$ so that $\eta_i|_{I_k}=0$ for every $k>i$.
    Since the added terms lie in $C$, the class of $\eta_i$ in $E=\row(\lambda)/C$ is still $\tau_i$.
        
    Now consider the $i$th factor.
    Every element of $D_i$ has $\varepsilon_i$-coefficient equal to $1$.
    Hence, for every adjacent pair, $d_{i,j}(\tau_i)=1$.
    Therefore the two rows $\chi_i|_{I_i}$ and $\eta_i|_{I_i}$ define a characteristic matrix on the polygon $P_{m_i}$.
    Indeed, on every adjacent pair the two column vectors are $\binom{1}{a}$ and $\binom{1}{a+1}$, which are linearly independent over $\Z_2$.
        
    Finally, consider the following ordered basis of $\row(\lambda)$
    $$
            \chi_1,\eta_1,\chi_2,\eta_2,\ldots,\chi_n,\eta_n.
    $$
    It is a basis because $\chi_1,\ldots,\chi_n$ form a basis of~$C$, and $\eta_1,\ldots,\eta_n$ project to the basis $\tau_1,\ldots,\tau_n$ of $E=\row(\lambda)/C$.
        
    Let $\lambda'$ be the matrix whose rows are ordered as above.
    Since $\lambda'$ and $\lambda$ have the same row space, they differ by left multiplication by an element of $GL(2n,\Z_2)$.
    Hence they are Davis--Januszkiewicz equivalent.
        
    By construction, the row pair $(\chi_i,\eta_i)$ has zero restriction to $I_k$ whenever $k>i$.
    Therefore $\lambda'$ is block lower triangular.
    Moreover, the $i$th diagonal block is
    $$
        \Lambda_{ii}
        =
        \begin{pmatrix}
        \chi_i|_{I_i}\\
        \eta_i|_{I_i}
        \end{pmatrix},
    $$
    which is a characteristic matrix on $P_{m_i}$.
    The first row of $\Lambda_{ii}$ is $\one_{m_i}$.
    This proves the theorem.
\end{proof}

\end{document}
